%% ===========================================================================
\FloatBarrier
\section{Numerical results and discussion}
\label{sec:results}
%% ===========================================================================

\FloatBarrier
\subsection{Material and simulation parameters}
\label{sec:material}

The material is tetragonal $\mathrm{BaTiO_3}$ at room temperature, with the
elastic, piezoelectric and dielectric constants measured by Zgonik and
co-workers~\cite{Zgonik1994}; Table~\ref{tab:params} collects them, the values
used in the runs and the status of each value. The dielectric anisotropy ratio
$\kappa_{11}/\kappa_{33}=39$ is retained, so that the two poling configurations
differ in their permittivity as well as in their modulus. The elastic,
piezoelectric and dielectric constants are measured, whereas the mass density
and the flexoelectric coefficient are \emph{nominal}: the density scales every
frequency by a common factor, and the flexoelectric coefficient enters every
bending coupling as $\mu^{2}$.



The strain-gradient length scale of the model is set by the window. With
half-widths $a_1\times a_2=3\times1$\,\si{\micro\metre} the cuboid window of
\eqref{eq:Mshapes} gives $M_{11}=3.0\times10^{-12}$ and
$M_{22}=3.3\times10^{-13}$\,\si{\square\metre}, that is characteristic lengths
$\sqrt{M_{11}}=\SI{1.73}{\micro\metre}$ and $\sqrt{M_{22}}=\SI{0.58}{\micro\metre}$
along and across the strip. The ratio $a_1/L=0.015$ is far inside the operating range
$a/L\le0.2$ that Sec.~\ref{sec:limit_window} assumes, so the second-order
truncation of \eqref{eq:Mdef} is not what limits the results; the strip is
solved with a $48\times4$ BFS mesh, which the convergence study of
Sec.~\ref{sec:verification} shows to be adequate for every quantity quoted
here.

\FloatBarrier
\subsection{Verification and convergence}
\label{sec:verification}

The numerical evidence is summarized here in a form that a reader can check
against a re-implementation; the complete set of benchmarks, identities and
independent routes is collected in \ref{app:verification}.

\emph{Benchmarks.} Three problems with known answers were solved with the same
element: a one-dimensional strain-gradient bar, whose static tip displacement
and frequency were compared with the closed form; the same bar in the classical
limit; and the constant-strain patch test. The clamped-free and clamped-clamped
frequencies of the strip agree with the Euler--Bernoulli values to $0.26\,\%$
and $0.10\,\%$, and the elastic strip reproduces the gradient theory whenever
the window is set to zero.

\emph{Mesh convergence.} Figure~\ref{fig:convergence} shows the error in the
frequency of the extensional mode and in its quasi-static coupling factor
against the number of degrees of freedom, for the strip of
Table~\ref{tab:params}. Both converge from above at the rate expected of a
bicubic Hermite interpolation, and the $48\times4$ mesh used throughout the
paper is inside the asymptotic range: refining it to $96\times4$ changes the
frequency by $0.05\,\%$ and the coupling factor by $0.03\,\%$, and the coarser
$24\times2$ mesh used for the parameter sweeps differs from it by $0.04\,\%$
and $0.03\,\%$. Table~\ref{tab:convergence} of \ref{app:data} lists the
four meshes. The plane-stress and plane-strain reductions are compared with the
three-dimensional solution in \ref{app:verification}, which is the one
place where the two-dimensional idealization is visible.

\emph{Time discretization.} The L1 scheme is first order in $\Delta t$ for
$\alpha<1$ and exact for $\alpha=1$. Halving the step at fixed
$\omega\tau=0.3$ changes the damping ratio extracted from the ring-down by less
than $0.5\,\%$ for $\alpha\ge0.5$, and the step used in Sec.~\ref{sec:transient}
is the halved one.

\emph{Closed-form coupling.} Table~\ref{tab:coupling} collects every coupling
factor quoted in this paper and the measure behind each. The first of them
comes from the constitutive tensors alone: the quasi-static coupling factor of
the electrode pair is available in closed form,
$k^2=(e^{\mathsf T}C^{-1}e/\kappa)/(1+e^{\mathsf T}C^{-1}e/\kappa)$, which for
the extensional configuration gives $x=e^{\mathsf T}C^{-1}e/\kappa=1.2768$ and
therefore $k^2=0.5608$ (Table~\ref{tab:closedform}). Three further closed
forms bracket it: the
three-dimensional free bar gives $0.5771$, the three-dimensional uniaxial case
$0.3748$ and the two-dimensional plane-strain case $0.5132$
(Table~\ref{tab:closedform}). The separation
between the plane-stress and plane-strain closed forms is a modeling choice,
and the plane-stress one is adopted because the strip is free out of plane.

The assembled route returns several numbers near $0.56$, and they are different
measures of the same tensors, so Table~\ref{tab:coupling} lists every coupling
factor quoted in this paper with its boundary condition, mesh, electrode pair,
matrix and formula. Three of them matter here. With the window switched off and
no constraint, the assembled route reproduces the closed form to
$1.2\times10^{-11}$ relative ($0.56079$), which is the validation of the
assembly. Clamping the left edge lowers it to $0.5596$ ($0.55957$) and
switching the window on lowers it further to $0.55814$: the first is the
constraint, the second the gradient stiffness, and together they explain the $0.36\,\%$ between
the closed form and the working value $0.5581$. No attempt is made to make the
three agree, and the abstract quotes the working value $0.5581$
($0.55814$ to five decimals) of the clamped $48\times4$ mesh with the window
on.

\begin{table}[!htbp]
\centering
\footnotesize
\caption{Every coupling factor quoted in this paper, one row each, with the
measure that produced it. All rows are plane stress, poled along $x$, with the
electrode pair on the short faces; the window is
$a_1\times a_2=3\times1$\,\si{\micro\metre} where it is switched on. The
mode-free rows are $x/(1+x)$ with $x=d^{\mathsf T}K_{sc}^{-1}d/\beta$ and use no
eigenvector; the mode-based rows are the extensional mode (mode~4) of the
$24\times2$ mesh, $(\omega_{oc}^{2}-\omega_{s}^{2})/\omega_{oc}^{2}$ with $s$
the matrix named in the fifth column. Rows 1 and 8 are the same number by
construction and are the validation of the assembled route; row 2 isolates the
clamped edge, row 3 the window, row 4 the mesh, and rows 6 and 7 differ only in
$\mu$.\ The last column is the entry of
\texttt{results/coupling\_provenance.json} that reproduces the row; that file
also records the upstream key each entry was checked against.}
\label{tab:coupling}
%% tabular body generated -- do not edit (see AUTHOR_TODO.md)
\setlength{\tabcolsep}{4pt}
\begin{tabular}{@{}r p{0.22\textwidth} l l p{0.26\textwidth}@{}}
\toprule
$k^{2}$ & build & mesh & $\mu$ & measure \\
\midrule
$0.56079$ & free, no window & 48$\times$4 & $0$ & $K_{sc}$, $x/(1+x)$ \\
$0.55957$ & left edge clamped, no window & 48$\times$4 & $0$ & $K_{sc}$, $x/(1+x)$ \\
$0.55814$ & left edge clamped, window on & 48$\times$4 & $0$ & $K_{sc}$, $x/(1+x)$ \\
$0.55797$ & left edge clamped, window on & 24$\times$2 & $0$ & $K_{sc}$, $x/(1+x)$ \\
$0.46498$ & left edge clamped, window on & 24$\times$2 & $10^{-6}$ & $K_{oc}$/$K_{uu}$, $k^{2}_{u}$ \\
$0.45830$ & left edge clamped, window on & 24$\times$2 & $10^{-6}$ & $K_{oc}$/$K_{sc}$, $k^{2}_{sc}$ \\
$0.45797$ & left edge clamped, window on & 24$\times$2 & $0$ & $K_{oc}$/$K_{sc}$, $k^{2}_{sc}$ \\
$0.56079$ & none (constitutive tensors) & --- & --- & $K_{sc}$, $x/(1+x)$ \\
\bottomrule
\end{tabular}

\end{table}

\emph{Independent routes.} Eight headline numbers were recomputed along routes
that do not use the assembled matrices at all -- the bar frequency from
Euler--Bernoulli theory, the coupling factor from the constitutive tensors, the
window coefficients from the geometry of \ref{app:Mcoeff}, and the
damping ratio from the analytical Kelvin--Voigt limit. All agree with the
finite element values within the tolerances stated in
\ref{app:verification}.

\begin{figure}[!htbp]
\centering
\includegraphics[width=\textwidth]{F2_convergence.pdf}
\caption{Convergence of the discretization, against the number of mechanical
degrees of freedom of four meshes. (a) Relative error of the extensional
coupling factor $k^2$ and of the frequency of the mode that carries it, both
referred to the finest mesh, with the $N^{-2}$ rate line expected of a bicubic
Hermite interpolation. Both are the mode-free quasi-static $k^2$ of the clamped
build with the window on; (b) the two quantities themselves, with the converged
value as a dotted line. The $48\times4$ mesh used in the paper is the rightmost
point of both panels and is marked.
($L=\SI{200}{\micro\metre}$, $h=\SI{10}{\micro\metre}$, $E$ and $\nu$ from
Table~\ref{tab:params}, $a_1\times a_2=3\times1$\,\si{\micro\metre},
$a_1/L=0.015$, $\mu=10^{-6}$\,\si{\coulomb\per\metre}, plane stress.)}
\label{fig:convergence}
\end{figure}

\begin{table}[!htbp]
\centering
\small
\caption{The coupling factor obtained from the constitutive tensors alone, for
the four kinematics, together with the coupling strength
$e^{\mathsf T}C^{-1}e/\kappa$ from which it follows. The record is that of
Table~\ref{tab:params}, in plane stress and in plane strain. The three finite
element rows below the rule are all mode-free values of \eqref{eq:kstatic},
plane stress, poling $x$, on the meshes named, with the window switched off;
Table~\ref{tab:coupling} sets them beside the same measure with the window on.}
\label{tab:closedform}
% table generated from the archived results
\begin{tabular}{@{}lcc@{}}
\toprule
& $x$ & $k^2$ \\
\midrule
3-D free bar & 1.3647 & 0.5771 \\
3-D uniaxial ($k_{33}^2$) & 0.5996 & 0.3748 \\
2-D plane stress, from the record & 1.2768 & \textbf{0.5608} \\
\midrule
FEM, plane stress, free, $48\times4$ & 1.2768 & \textbf{0.56079} \\
FEM, plane stress, clamped, $24\times2$ & 1.2668 & 0.55884 \\
FEM, plane stress, clamped, $48\times4$ & 1.2705 & 0.55957 \\
FEM, plane stress, clamped, $96\times4$ & 1.2723 & 0.55992 \\
\midrule
2-D plane strain, from the record & 1.0541 & 0.5132 \\
FEM, plane strain, clamped, $48 \times 4$ & 1.0471 & 0.5115 \\
\bottomrule
\end{tabular}

\end{table}

Figure~\ref{fig:fields} records the four fields of the coupled bending mode in
the open-circuit configuration. The in-plane and the transverse displacement
are those of a first bending mode. The potential, by contrast, is not: it
reaches tens of millivolts a few tens of micrometres from the clamped edge
and near mid-height, and it falls toward both the edge and the free end, so
the two face electrodes sit where the potential is smallest. That is the
spatial form of the cancellation that makes this configuration insensitive to
the piezoelectric charge and responsive to the strain gradient
(Sec.~\ref{sec:thickness}).

\begin{figure}[!htbp]
\centering
\includegraphics[width=\textwidth]{F3_fields.pdf}
\caption{Fields of the coupled bending mode of the strip in the open-circuit
configuration, on the $48\times4$ mesh: (a) in-plane displacement $u_x$
normalized by its maximum, (b) transverse displacement $w$, (c) potential
$\phi$ at \SI{10}{\nano\metre} of tip motion, (d) magnitude of the electric
field. The potential peaks a few tens of micrometres from the clamped edge and
near mid-height, and is smallest at the face electrodes, which is why this
electrode pair collects the gradient rather than the strain. ($L=\SI{200}{\micro\metre}$,
$h=\SI{10}{\micro\metre}$, poling $y$, plane stress,
$a_1\times a_2=3\times1$\,\si{\micro\metre}, $\mu=10^{-6}$\,\si{\coulomb\per\metre}.)}
\label{fig:fields}
\end{figure}

\FloatBarrier
\subsection{Worked calculation of the headline numbers}
\label{sec:worked}

The numbers quoted in the abstract are reproduced here from the inputs, so that
the route from a material constant to a coupling factor is visible and not only
its outcome. Each value below is the one carried into the figures and tables of
this section.

\paragraph{Window.} For half-widths $a_1=\SI{3}{\micro\metre}$ and
$a_2=\SI{1}{\micro\metre}$ the second-moment tensor of the cuboid window is
diagonal with
\begin{equation}
  M_{11}=\frac{a_1^{2}}{3}=\SI{3.000e-12}{\metre\squared},\qquad
  M_{22}=\frac{a_2^{2}}{3}=\SI{3.333e-13}{\metre\squared},
  \label{eq:Mnumber}
\end{equation}
that is $M_{11}/a_1^{2}=M_{22}/a_2^{2}=1/3$ and an anisotropy ratio
$M_{11}/M_{22}=9$. The three shapes are interchangeable once $M$ is fixed,
because $M_{ii}=a_i^{2}/3$, $a_i^{2}/4$ and $a_i^{2}/5$ for the cuboid, the
ellipse and the ellipsoid: an ellipsoid of half-width
$a_{\rm ell}=1.2910\,a_{\rm cub}=\sqrt{5/3}\,a_{\rm cub}$ therefore reproduces
the cuboid value of \eqref{eq:Mnumber} exactly, and the reciprocal statement,
$a_{\rm cub}=0.7746\,a_{\rm ell}=\sqrt{3/5}\,a_{\rm ell}$, is the one used for
the equal-$M$ pair of Sec.~\ref{sec:windowshape}.

\paragraph{Plane-stress reduction.} The strip is free out of plane, so the
stiffness that multiplies $\varepsilon_{xx}$ is not $C_{11}$ but the reduced
entry
\begin{equation}
  C^{p}_{ij}=C_{ij}-\frac{C_{i3}C_{j3}}{C_{33}},
  \label{eq:Cp}
\end{equation}
which for $\mathrm{BaTiO_3}$ gives $C^{p}_{xx}=\SI{9.5500e10}{\pascal}$ with the
poling along $x$ and $\SI{1.6946e11}{\pascal}$ with the poling along $y$, with
$C^{p}_{12}=\SI{5.7000e10}{\pascal}$ in both cases. The entries that pair with
it are $e_{33}=6.7$ and $e_{31}=-0.7$\,\si{\coulomb\per\square\metre} and the
permittivity along the electrode gap, $\kappa_{yy}=\kappa_{33}=$
\SI{4.9583e-10}{\farad\per\metre}. The exchange of the two moduli between the
two poling directions, and not a change of material, is what makes the
extensional and the bending configuration different electrical problems
(Table~\ref{tab:axes}).

\paragraph{Closed form and the assembled route.} From the quantities above,
\begin{equation}
  x=\frac{e^{\mathsf T}C_{p}^{-1}e}{\kappa}=1.2768,
  \qquad
  k^{2}=\frac{x}{1+x}=0.5608,
  \label{eq:xclosed}
\end{equation}
which involves no mesh. The same quantity on the assembled model is obtained
from one short-circuit and one open-circuit reduction of the $48\times4$ mesh:
$49\times5=245$ nodes, eight mechanical degrees of freedom per node
($1960$ mechanical unknowns) and $245$ potential unknowns, condensed to a
single electrode pair $(d,\beta)$ with $\beta=\SI{2.4792e-11}{\farad}$ and
$d$ the charge vector of that pair. Then
\begin{equation}
  x_{\rm FEM}=\frac{d^{\mathsf T}K_{sc}^{-1}d}{\beta}=1.2631,
  \qquad
  k^{2}_{\rm FEM}=\frac{x_{\rm FEM}}{1+x_{\rm FEM}}=0.5581,
  \label{eq:xFEM}
\end{equation}
$0.47\,\%$ below \eqref{eq:xclosed}. That residual is not discretization: it
is the clamped edge, which costs $0.22\,\%$, and the gradient stiffness of the
window, which costs a further $0.26\,\%$ (Table~\ref{tab:coupling}), while the
unconstrained strip without the window reproduces \eqref{eq:xclosed} to
$1.2\times10^{-11}$ relative and the mesh refinement of the working build moves
the value by $0.03\,\%$ (Table~\ref{tab:convergence}).

\paragraph{Flexoelectric bending.} With the piezoelectric constants switched
off, the same reduction gives the smallness parameter of the bending
configuration,
\begin{equation}
  x=\frac{3\mu^{2}}{C^{p}_{xx}\,h^{2}\,\kappa_{yy}}
   =\frac{3\,(10^{-6})^{2}}{1.6946\times10^{11}\cdot(10^{-5})^{2}\cdot
     4.9583\times10^{-10}}
   =\num{3.5704e-4}
  \label{eq:xflexo}
\end{equation}
at $h=\SI{10}{\micro\metre}$, against a finite element value of
\num{1.7437e-4}: the analytic route gives the order of the effect and the
factor by which it grows, the finite element gives its value. Both are listed in
Table~\ref{tab:worked}, where the local log--log slope of the finite element
series is formed as $\ln(k^2_{i+1}/k^2_{i})/\ln(h_{i+1}/h_i)$ on consecutive
thicknesses: $-1.96$ between $10$ and \SI{7}{\micro\metre}, $-1.52$ between $5$
and \SI{3.5}{\micro\metre}, and $-0.43$ between $1.4$ and
\SI{1}{\micro\metre}.

\paragraph{Memory.} The L1 weights of \eqref{eq:l1} are fixed numbers that do
not depend on the model,
\begin{equation}
  \alpha=0.5:\; (b_0,b_1,b_2)=(1,\;0.414214,\;0.317837),\qquad
  \alpha=0.7:\; (b_0,b_1,b_2)=(1,\;0.231144,\;0.159245),
  \label{eq:bweights}
\end{equation}
and the time march advances the modal state with the matrix
$M+\beta\Delta t^{2-\alpha}K_{\rm mem}/\Gamma(2-\alpha)$ of
\eqref{eq:stepA}, where $K_{\rm mem}=\tau^{\alpha}K$ and the stiffness $K$ is
either $K_{\rm eff}$ or $K_{uu}$. At $\omega\tau=0.3$ the ratio of two
successive voltage peaks is $2.5947$, so the decrement is $\delta=0.95345$ and
\begin{equation}
  \zeta_{\rm eff}=\frac{\delta}{\sqrt{4\pi^{2}+\delta^{2}}}=0.15003
  \quad\text{against}\quad
  \frac{\omega\tau}{2}=0.15000,
  \label{eq:zetaextract}
\end{equation}
an error of $0.02\,\%$; the same extraction on $K_{uu}$ gives $0.06746$, a
ratio $K_{uu}/K_{\rm eff}=0.450$.

\begin{table}[!htbp]
\centering
\small
\caption{Inputs, intermediate quantities and outputs for the headline numbers
of the abstract. Every entry is the value carried into the figures and tables
of this paper; the last column is the number the abstract quotes.}
\label{tab:worked}
\begin{tabular}{@{}p{0.24\textwidth}p{0.20\textwidth}p{0.26\textwidth}p{0.18\textwidth}@{}}
\toprule
output & inputs & intermediate quantity & value \\
\midrule
extensional $k^{2}$ & $C_{p}$, $e$, $\kappa_{33}$ & $x=1.2768$
(eq.~\ref{eq:xclosed}); $x_{\rm FEM}=1.2631$ (eq.~\ref{eq:xFEM}) &
$0.5581$ ($0.5608$ closed form) \\
thickness independence & $k^{2}$ at $10$ and \SI{1}{\micro\metre} &
$0.55814\to0.55861$ & $0.09\,\%$ \\
flexoelectric growth & $\mu$, $C^{p}_{xx}$, $h$, $\kappa_{yy}$ &
$x=3.5704\times10^{-4}$ at \SI{10}{\micro\metre};
$1.7437\times10^{-4}\to4.1736\times10^{-3}$ & factor $24$ \\
flexoelectric slope & the eight thicknesses & local log--log slopes &
$-1.96\to-0.43$ \\
window anisotropy & $M_{11}$, $M_{22}$, $\theta$ & $M$ rotates, the
mesh does not & $15$--$19\,\%$ \\
two configurations & \eqref{eq:xclosed} in both & $0.4580$ against $0.1141$;
\SI{138.99}{\volt} against \SI{2.41}{\volt}, equal tip displacement & factor $58$ \\
memory ratio & $K_{\rm eff}$, $K_{uu}$ & $\zeta=0.15003$ against $0.06746$ &
$0.450$ \\
integer-order limit & $\omega\tau=0.3$ & $\zeta_{\rm eff}$ against
$\omega\tau/2$ & $0.02\,\%$ \\
\bottomrule
\end{tabular}
\end{table}

\FloatBarrier
\subsection{Window orientation and strip aspect ratio}
\label{sec:orientation}

Figure~\ref{fig:orientation} varies the orientation $\theta$ of the rectangular
averaging window and the strip aspect ratio, for the poling-$x$
(length-extensional) configuration that dominates the response. The window is
held at a half-width of \SI{3}{\micro\metre} along its long axis; the strip is
\SI{200}{\micro\metre}$\times$\SI{10}{\micro\metre}. The extensional mode is
the one carrying the largest coupling, and over the $6\times19$ cuboid cases
that identification never changes, its short-circuit frequency varying smoothly
from $4907$ to \SI{5296}{kHz}.

Only the first quadrant is computed. The strip maps onto itself under
$\theta\to\theta+180^\circ$. The second invariance is a symmetry of the whole
problem and not of the energy alone: the gradient energy is \emph{linear} in
$M$, so a change of sign of $M_{12}$ by itself changes the response, and what
leaves it unchanged is the reflection $y\to-y$, which maps $\theta\to-\theta$
and at the same time maps the strip, the electrode pair and -- for poling
along $x$, the configuration plotted here -- the material onto themselves, the
$4mm$ poling axis lying in the plane along $x$. The response is therefore even
in $\theta$ for the poling-$x$ configuration. For poling along $y$ the same
reflection reverses the poling direction and is not a symmetry of the material,
so evenness in $\theta$ does not follow there and only the period
$\theta\to\theta+180^\circ$ holds; the orientation study is confined to
poling $x$ for that reason. Neither invariance is assumed: both are evaluated
at \SI{15}{\degree} intervals for the extreme aspect ratios, where they
reproduce the computed points to \num{3e-11} relative. The full circles in
Figs.~\ref{fig:orientation}(a) and~(b) are therefore the computed quadrant
mirrored twice and joined by a periodic cubic spline, the markers are the
computed points, and the maps in (c) and (d) are drawn over
$0$--$180^\circ$ for the same reason.

Both quantities move monotonically with $\theta$ and both saturate by
$\theta=90^\circ$. The extensional frequency rises from $5012$ to
\SI{5296}{kHz} at ${\rm AR}=2$ and from $4907$ to \SI{5288}{kHz} at
${\rm AR}=10$, that is by 6--8\,\%, while the coupling falls from $0.4448$ to
$0.3786$ and from $0.4677$ to $0.3801$ between the same two orientations, a
change of 15--19\,\%. The fundamental mode of the same strip, which carries a
coupling three orders of magnitude smaller, moves from \SI{146.4}{kHz} to
\SI{158.1}{kHz} at ${\rm AR}=10$; the two modes are not interchangeable, and
they are named explicitly wherever both appear. At $\theta=0^\circ$ the window
aspect ratio is worth $0.4448\to0.4677$ between ${\rm AR}=2$ and ${\rm AR}=10$,
about 5\,\%, and the aspect ratios collapse onto $0.3786$--$0.3801$ at
$\theta=90^\circ$: once the long axis of the window is transverse to the strain
gradient, its aspect ratio no longer matters.

Two remarks keep this in proportion. First, the coupling plotted here is the
\emph{total} coupling of the extensional mode and it is piezoelectric in the
main; the flexoelectric contribution at $\mu=10^{-6}$\,\si{\coulomb\per\metre}
is of order $10^{-4}$ (Sec.~\ref{sec:thickness}), so the window acts on the
modal coupling through the gradient stiffness in $K_{sc}$ rather than by adding
a flexoelectric term directly. Second, in the language of
Sec.~\ref{sec:window} the window enters only through $M$: rotating the window
rotates $M$, and the 15--19\,\% spread is a property of the second-moment
tensor of a rectangular window, not of the rectangle as such.
Figure~\ref{fig:design} separates the two. Panel (a) is the change of the
coupling against the aligned window: largest for the most elongated window,
$-18.7$\,\% at ${\rm AR}=10$ and $\theta=90^\circ$, and weakest for the
squarest, $-14.9$\,\% at ${\rm AR}=2$. Panel (b) is
$(M_{11}-M_{22})/(M_{11}+M_{22})$, a function of the orientation alone, which
passes through zero at $45^\circ$ and changes fastest there: the coupling
follows the anisotropy of the tensor and not the shape of the window.

\begin{figure}[!htbp]
\centering
\includegraphics[width=\textwidth]{F4_orientation.pdf}
\caption{Window orientation at fixed window size, on
\SI{200}{\micro\metre}$\times$\SI{10}{\micro\metre} strips with
$h=\SI{10}{\micro\metre}$, poling $x$, plane stress, $24\times2$ mesh, $\mu=0$
and $\alpha=1$. (a) Frequency of the extensional mode and (b) its coupling
factor, both normalized to the aligned window, over the full circle; the
radial axis is truncated and the dotted circle marks unity. Markers are the
computed points, $5^\circ$ apart in the quadrant and mirrored; the lines are a
periodic cubic spline through them. (c), (d) The same two quantities as maps
over the window angle and the aspect ratio, on $0$--$180^\circ$ because the
response repeats under $\theta\to\theta+180^\circ$; white dots mark the six
computed aspect ratios. The polar sections and the maps are built from the
same $19\times6$ cuboid sweep, plus one elliptical window of equal $M$ in (b).}
\label{fig:orientation}
\end{figure}

\begin{figure}[!htbp]
\centering
\includegraphics[width=0.86\textwidth]{F5_design.pdf}
\caption{The same sweep as Fig.~\ref{fig:orientation}, read as a design space.
(a) Change of the extensional coupling relative to the aligned window, one
curve per aspect ratio, with markers at every second computed orientation; the
curves span $4$ percentage points at $90^\circ$ and the deepest is
$-18.7$\,\%. (b) Anisotropy of the second-moment tensor,
$(M_{11}-M_{22})/(M_{11}+M_{22})$, which is independent of the aspect ratio and
of the strip, and whose zero at $45^\circ$ is where the coupling in (a)
changes fastest. Same strip and mesh as Fig.~\ref{fig:orientation}.}
\label{fig:design}
\end{figure}

\FloatBarrier
\subsection{Fractional memory and the damping map}
\label{sec:memorystudy}

Two questions are separated here. The first is \emph{model}: the fractional
Kelvin--Voigt element of \eqref{eq:kv} is one way to represent dielectric loss,
and it differs measurably from the integer-order element that finite element
codes usually provide. The second is \emph{implementation}: the convolution of \eqref{eq:step} may multiply the
fully condensed stiffness $K_{\rm eff}$ or only the short-circuit matrix
$K_{uu}$, and the two give different damping.

\paragraph{Memory matrix.} Table~\ref{tab:memory} reports the ring-down of the
extensional mode (mode 4, \SI{6.727}{MHz}) at
$\tau=\SI{7.0974}{\nano\second}$, that is $\omega\tau=0.3$, with an initial tip
displacement that produces \SI{139.07}{\volt} of open-circuit voltage at
$t=0$. Three coupling numbers for that mode appear in the paper. They come
from two quantities, not three: $k^2_{4}(K_{uu})=0.4650$ is the mode-based
measure on the elastic memory matrix $K_{uu}$ (Table~\ref{tab:memory}) and
$k^2_{4}(K_{sc})=0.4583$ is the same measure on the short-circuit stiffness
$K_{sc}$, both on the $24\times2$ mesh of the ring-down; the third, $0.4580$ in
Sec.~\ref{sec:application}, is $k^2_{4}(K_{sc})$ again, recomputed with
$\mu=0$ because that section quotes the material without flexoelectricity.
Table~\ref{tab:coupling} lists all three. The gap of $1.5\,\%$ between
$0.4650$ and $0.4583$ is therefore the choice of memory matrix and nothing
else, and the $0.07\,\%$ between $0.4583$ and $0.4580$ is $\mu$; neither pair
is averaged, and the ring-down below uses $K_{uu}$ and $K_{\rm eff}$. The damping ratio is extracted from two successive peaks through the exact
relation for an exponentially decaying oscillation,
\begin{equation}
  \delta = \ln\frac{A_k}{A_{k+1}},\qquad
  \zeta_{\rm eff} = \frac{\delta}{\sqrt{4\pi^{2}+\delta^{2}}},
  \label{eq:zeta}
\end{equation}
of which $\zeta=\delta/2\pi$ is the small-damping approximation; at the values
below the two differ by $1.2\,\%$, so the approximation is not used. At the
$\omega\tau=0.3$ of Table~\ref{tab:memory} the ratio of two successive voltage
peaks is $2.5947$, so $\delta=0.95345$ and $\zeta_{\rm eff}=0.15003$ for
$K_{\rm eff}$ at $\alpha=1$, against the exact $\omega\tau/2=0.15000$: the
implementation reproduces the classical limit to $0.02\,\%$. Refining the step,
$\zeta_{\rm eff}=0.15011$, $0.15006$ and $0.15003$ at $N$, $2N$ and $4N$ steps
over twenty periods, so the residual is first order in $\Delta t$ and
extrapolates to $0.15000$; the order is the one expected of the L1 weights at
$\alpha=1$, whose right-hand end point makes the scheme behave as backward
Euler. Replacing $K_{\rm eff}$ by $K_{uu}$ removes the piezoelectric
contribution from the memory and takes the ratio
$\zeta(K_{uu})/\zeta(K_{\rm eff})$ at $\alpha=1$ down to $0.450$, so the
short-circuit memory dissipates $55\,\%$ less. The difference is of the order
of the coupling ($k^2_{4}=0.4650$) but does \emph{not} reduce to the
naive factor $1-k^2=0.535$: the two memories differ by the piezoelectric
stiffness that the condensation puts into $K_{\rm eff}$, and its effect on the
ring-down has to be computed rather than estimated. For a strongly
piezoelectric device the distinction therefore belongs in the model statement,
not in the implementation.

\begin{table}[!htbp]
\centering
\small
\caption{Ring-down of the extensional mode (mode 4, \SI{6.727}{MHz},
$k^2_{4}(K_{uu})=0.4650$, $\omega\tau=0.3$,
$\tau=\SI{7.0974}{\nano\second}$), for memory carried on the condensed
stiffness $K_{\rm eff}$ and on the short-circuit stiffness $K_{uu}$. At
$\alpha=1$ the entry is the damping ratio \eqref{eq:zeta} of the exponential
decay. At $\alpha=0.7$ the free decay of a fractional element is not
exponential, so no single damping ratio exists: the entry there is the
first-period local ratio of the first two peaks, and the successive decrements
are listed in the text. The last row is the ratio of the two entries.}
\label{tab:memory}
\begin{tabular}{@{}lcc@{}}
\toprule
memory matrix & $\alpha=1$: $\zeta_{\rm eff}$ &
  $\alpha=0.7$: first-period $\zeta$ \\
\midrule
$K_{\rm eff}$ & $0.15003$ & $0.1942$ \\
$K_{uu}$      & $0.06746$ & $0.0831$ \\
\midrule
ratio $K_{uu}/K_{\rm eff}$ & $0.450$ & $0.428$ \\
\bottomrule
\end{tabular}
\end{table}

\paragraph{What the two memories stand for.} Neither matrix is derived from a
dielectric relaxation; both are placements of the fractional element inside a
condensed model, and the choice is treated here as a model statement rather
than as a material property. One law can be derived exactly. If the
permittivity itself relaxes,
$D_k=e_{kij}\varepsilon_{ij}+\kappa_{kl}[1+(\tau\partial_t)^{\alpha}]E_l$,
then $\nabla\cdot\bm D=0$ carries the Caputo derivative of the field, and
condensing the potential moves the memory onto the piezoelectric stiffening
$B_iC^{-1}B_i^{\mathsf T}$, that is onto exactly the part of $K_{\rm eff}$ that
the condensation adds over $K_{uu}$. That is a third possibility, memory on
the electromechanical block alone, bracketed by the two matrices tabulated:
$K_{\rm eff}$ remembers the whole condensed stiffness and $K_{uu}$ none of it.
The ratio $0.450$ is a statement about the $\alpha=1$ row of
Table~\ref{tab:memory} and about this device, not a universal factor.

\paragraph{Decay at $\alpha<1$.} At $\alpha=0.7$ the free decay is not
exponential and a single damping ratio does not exist: the successive
decrements of the $K_{\rm eff}$ ring-down are $1.244$, $1.547$, $0.010$,
$0.403$, $0.240$ and $0.151$; the first, second, fourth, fifth and sixth of
them are unchanged to three digits between $\Delta t$, $\Delta t/2$ and
$\Delta t/4$, so they are a property of the fractional element -- whose
relaxation is a Mittag-Leffler function, not an exponential -- and not of the
step, while the third is a near-cancellation between two peaks of comparable
height, where a ratio of amplitudes does not resolve a decrement. A straight line fitted to $\ln A_k$ over
the twenty periods therefore returns a number that is not a damping ratio,
which is why the $\alpha=0.7$ column of Table~\ref{tab:memory} carries the
first-period decrement rather than a fitted slope.
The $\alpha=1$ entries are unaffected: there the twenty successive decrements
agree to $7\times10^{-10}$ relative, which is why \eqref{eq:zeta} applies.

\paragraph{Fractional against integer-order element.} Table~\ref{tab:modulus}
compares the complex modulus of the fractional element \eqref{eq:kvfreq} with
the integer-order Kelvin--Voigt element and a Zener element fitted to it at the
device point $\omega\tau=0.3$ and held fixed afterwards. At that point the
Kelvin--Voigt model over-predicts the modulus by $3.5\,\%$ at $\alpha=1$ and by
$22.6\,\%$ at $\alpha=0.7$; at $\omega\tau=1$ the ratio of the two is $1.35$
and $1.62$, and at $\omega\tau=100$ the integer-order model is wrong by a
factor $91$. A
single-relaxation fit is unusable outside the narrow band where it was fitted,
which is the band a design map is meant to cover.



\paragraph{The damping map.} Figure~\ref{fig:damping}(a) is the pole of
\eqref{eq:kvfreq} in the complex plane, traced over $\omega\tau$ for seven
relaxation times and coloured by $\tau$, with markers at five Caputo orders;
Fig.~\ref{fig:damping}(b) is the damping ratio
$\zeta=-\mathrm{Re}\,s/(\omega_0|s|)$ over the $(\alpha,\omega\tau)$ plane.
Both are computed with branch continuation: the root is followed from one
$\omega\tau$ to the next by warm-starting the Newton iteration at the previous
pole, without which the iteration jumps to a spurious root of the
negative-real-part branch.

All poles lie in the left half-plane for every $(\alpha,\tau)$ examined, so the
fractional Kelvin--Voigt element is dissipative by construction. The
integer-order case ($\alpha=1$) sits at $|\mathrm{Im}\,s|/\omega_0\simeq1$ for
$\omega\tau\lesssim0.1$ and moves down and to the right as $\omega\tau$ grows,
so the dashpot lowers the frequency and eventually over-damps the mode: at
$\omega\tau=10$ the imaginary part has vanished and $\zeta=1$. The fractional
cases ($\alpha<1$) start from the same neighborhood and fan \emph{upward}: a
spring-pot with $\alpha<1$ raises the damped frequency as $\tau$ grows, because
the fractional modulus grows sub-linearly rather than linearly in $\omega$. At
$\omega\tau=10$ and $\alpha=0.5$ the damped frequency is $2.13\,\omega_0$
($2.45\,\omega_0$ at $\alpha=0.7$), against $\mathrm{Im}\,s=0$ for $\alpha=1$
at the same $\omega\tau$. A single-fit Kelvin--Voigt parameterization cannot
reproduce that, and the pole rather than the modulus is what a design should be
read from. The ratio in Fig.~\ref{fig:damping}(b) increases with both $\alpha$
and $\omega\tau$ and reaches $\zeta\to1$ for $\alpha\to1$ and large
$\omega\tau$, where the element is critically damped; for small $\alpha$ it is
flat, because a spring-pot with a broad relaxation spectrum produces a damping
that depends only weakly on the device frequency. That flatness is the
qualitative reason the fractional description is used for dielectric loss in
the first place. The band $\omega\tau\le1$, which contains the strip mode at
\SI{6.73}{MHz} for every $\tau\le\SI{32}{nano\second}$ used here, is the one in
which the element is calibrated and claimed; the wider range in panel (b) is
drawn for continuity rather than as a prediction.

\begin{figure}[!htbp]
\centering
\includegraphics[width=\textwidth]{F6_damping.pdf}
\caption{The fractional Kelvin--Voigt element in the complex plane and as a
damping map, from the pole of \eqref{eq:kvfreq} with branch continuation.
(a) The pole of a mode with undamped frequency $\omega_0=\SI{6.73}{MHz}$,
traced over $\omega\tau=0.02\ldots1.35$ for seven relaxation times
$\tau=0.5\ldots\SI{32}{nano\second}$ and coloured by $\tau$; markers are
$\alpha=0.2,\,0.4,\,0.6,\,0.8,\,1$, and the dotted line is
$\mathrm{Im}\,s=\omega_0$. Every pole has a negative real part, and the
fractional curves fan upward because the spring-pot modulus grows
sub-linearly in frequency. (b) Damping ratio
$\zeta=-\mathrm{Re}\,s/(\omega_0|s|)$ over the wider plane
$\omega\tau=10^{-1.5}\ldots10^{1.5}$ and $\alpha=0.2\ldots1$, with the same
branch continuation; contours every $0.1$. The ratio here is that of the
complex pole, $\zeta=-\mathrm{Re}\,s/(\omega_0|s|)$, and not a decrement of a
time record: for $\alpha<1$ the two coincide only to the extent that the free
decay is exponential, so the ring-down of
Sec.~\ref{sec:transient} is read with the first-period decrement of its table
and this pole ratio is used as a map of the element rather than as a
prediction of the decay.}
\label{fig:damping}
\end{figure}

\FloatBarrier
\subsection{Window shape versus second moment}
\label{sec:windowshape}

Figure~\ref{fig:window}(a) is the numerical statement of
Sec.~\ref{sec:window} on an isotropic plane-strain test problem (rectangle
$1\times0.2$, $20\times4$ mesh, uniaxial edge traction, displacement $u_R$
averaged over the loaded edge), plotted against the normalized window size for
an ellipsoidal window of semi-axis $l$, a cuboidal window of half-width $l$,
and a cuboidal window scaled to carry the same $M$ as the ellipsoid,
$a=l\sqrt{3/5}$. The third curve is indistinguishable from the first: within the adopted
second-moment model the shape of the window is irrelevant at fixed $M$, and
only $M$ enters it.
The second curve shows the error made by ignoring the distinction: equal
half-widths instead of equal $M$ over-stiffen the response, and the
discrepancy grows with the window size. At $l/L_x=0.2$ a cuboid of half-width
$l$ retains $0.922$ of the classical loaded-edge displacement, an ellipsoid of
semi-axis $l$ retains $0.941$, and the equal-$M$ cuboid sits on the ellipsoid
curve to the third decimal. At fixed $M$ the shape of the window is therefore
immaterial to $10^{-3}$, which is the coefficient ratio $5/3$ in the gradient
stiffness seen from the response side. The range in which the model is used is
$a/L\le0.2$, the right-hand edge of this figure. No larger window is run, and
the reason is that beyond this point we can no longer check the second-moment
truncation itself, not that we have located where it fails
(Sec.~\ref{sec:limit_window}).

Panels (b) and (c) are the corresponding static statement. The clamped strip of
Table~\ref{tab:params} is solved under an end load for three window sizes and
each solution is compared with the classical one of the same geometry and load.
The gradient terms stiffen the member. The difference against the classical
solution vanishes at the clamped edge, where both fields satisfy the same
kinematic conditions, and then accumulates along the member instead of decaying
over a thin layer: at these window sizes the gradient effect is distributed
rather than local, and its size grows with the window. That is why
$a/L\le0.2$ is quoted in Sec.~\ref{sec:limit_window} as an assumed operating
range and not as a validated truncation threshold: it is where the model is
used, and it is not a mesh requirement.

\begin{figure}[!htbp]
\centering
\includegraphics[width=\textwidth]{F7_window.pdf}
\caption{What the window size and shape do, on problems with an exact
classical counterpart. (a) Loaded-edge displacement of an isotropic
plane-strain rectangle ($E$ and $\nu$ of Table~\ref{tab:params}, $a/L_x\le0.2$,
$20\times4$ mesh, uniaxial edge traction) against the normalized window size,
for an ellipsoidal window of semi-axis $l$, a cuboidal window of half-width
$l$, and a cuboidal window scaled to the same second moment, $a=l\sqrt{3/5}$;
the response is normalized by its classical value, and the values at
$l/L_x=0.2$ are printed. (b) Deflected shape of the clamped strip under an end
load, normalized by the tip value, for three window sizes
($L=\SI{200}{\micro\metre}$, $h=\SI{10}{\micro\metre}$, $48\times4$ mesh,
poling $y$, plane stress, $\mu=10^{-6}$\,\si{\coulomb\per\metre},
$a_2=\SI{1}{\micro\metre}$). (c) Difference from the classical profile of the
same problem, normalized by the tip deflection: it vanishes at the clamped edge
and accumulates along the member rather than decaying over a boundary layer, at
all three window sizes. Markers are computed points, lines are interpolations.}
\label{fig:window}
\end{figure}

\FloatBarrier
\subsection{Thickness: a flexoelectric size effect against a
thickness-independent piezoelectric coupling}
\label{sec:thickness}

Figure~\ref{fig:thickness} is the central design result. The same strip is
solved at eight thicknesses between \SI{10}{\micro\metre} and
\SI{1}{\micro\metre} with a \SI{200}{\micro\metre} length and the window held
at half-widths $3\times1$\,\si{\micro\metre}, so that the thickness is the only
variable. One qualification belongs to the two thinnest points: the window
half-thickness is \SI{1}{\micro\metre} throughout, so at $h=\SI{2}{\micro
\metre}$ it is half the strip and at $h=\SI{1}{\micro\metre}$ it equals the
strip. The second-moment tensor of the full window is used at every thickness
and no truncated-window boundary correction is applied, so $M$ acts there as a
local effective gradient tensor rather than as the moment of a window small
compared with the body; the results at $h\le\SI{2}{\micro\metre}$ are read
under that assumption, which is consistent with the operating range $a/L\le0.2$ used
elsewhere but not verified independently at these thicknesses. The panels are
two electrical configurations of that strip, not two
estimates of one number: panel~(a) is poled across the thickness (electrodes on
the long faces, gap $h$, $\kappa_{yy}=\kappa_{33}$) and panel~(b) along the
length (electrodes on the short faces, gap $L$, $\kappa_{xx}=\kappa_{33}$), as
set out in Table~\ref{tab:axes}. Panel~(a) is the flexoelectric bending
coupling with $\mu=10^{-6}$\,\si{\coulomb\per\metre}. Panel~(b) is the
piezoelectric length-extensional coupling of the same strip, evaluated with
$\mu=0$ and with the mode-free measure \eqref{eq:kstatic}.

\begin{figure}[!htbp]
\centering
\includegraphics[width=\textwidth]{F8_thickness.pdf}
\caption{(a) Flexoelectric bending coupling versus strip thickness
($L=\SI{200}{\micro\metre}$, window half-widths
$3\times1$\,\si{\micro\metre}, $\mu=10^{-6}$\,\si{\coulomb\per\metre}, poling
$y$, plane stress) at eight thicknesses, with three of the local log--log
slopes annotated, a $96\times4$ witness at the thickest and the thinnest point
(open squares, indistinguishable from $48\times4$), and the $h^{-2}$ guide line
through the thinnest point. (b) Piezoelectric length-extensional coupling of
the same strip with $\mu=0$, mode-free: constant to $0.09\,\%$ over
$10\to1$\,\si{\micro\metre} and $0.4\,\%$ below the constitutive closed form
$0.5608$ (dotted), which is the discretization error of the $48\times4$ mesh
rather than a modeling difference. (c) Ratio of the two couplings on a
logarithmic scale: it spans $4.3\times10^{-4}$ to $7.5\times10^{-3}$ over the
range, so the flexoelectric term is not competitive at these parameters and the
two effects are not additive. ($\alpha=1$, $a_1/L=0.015$,
$L=\SI{200}{\micro\metre}$.)}
\label{fig:thickness}
\end{figure}

\paragraph{Flexoelectric size effect.} The flexoelectric coupling rises by a
factor $24$ from \num{1.74e-4} at \SI{10}{\micro\metre} to \num{4.2e-3} at
\SI{1}{\micro\metre}. Its local log--log slope is $-1.96$ at the thick end,
close to the $h^{-2}$ scaling of a gradient effect with a window much smaller
than the thickness, and falls to $-1.52$, $-1.26$ and $-0.43$ on the
successively thinner pairs, the seven segments spanning $-1.96$ to $-0.43$. The
saturation is not a mesh effect: a $96\times4$ witness reproduces the
$48\times4$ values to $0.04\,\%$ and $0.35\,\%$ at the two ends of the range,
and the residual at $\mu=0$ is $1.0\times10^{-10}$, seven orders of magnitude
below the signal. The dimensionless group that could explain the saturation is
$a_2/h$, the window half-thickness over the strip thickness, which runs from
$0.10$ at $h=\SI{10}{\micro\metre}$ to $1.0$ at \SI{1}{\micro\metre} on the
same seven segments: the slope decreases monotonically as $a_2/h$ grows, which
is consistent with a fixed-size window smearing the through-thickness strain
gradient once it becomes comparable to the strip. That reading is a hypothesis
of this study and is not tested here independently of the seven slopes above;
the same stiffening appears in the fundamental frequency, which follows
$210.4,108.0,50.2,34.7$\,kHz where pure Euler--Bernoulli scaling from the
thickest strip would give $210,105,42,21$\,kHz.

Figure~\ref{fig:modetrack} records the mode identification. The mode shape
computed at each thickness is compared with the one at
$h=\SI{10}{\micro\metre}$ through the modal assurance criterion, which stays
above $0.999$ across the whole sweep, and the ratio $f_2/f_1$ stays between
$6.21$ and $6.27$: the branch that carries the coupling is separated from its
neighbours by the same margin at every thickness, so continuity of the shape
and not an eigenvalue threshold selects it in every result of this paper.

\begin{figure}[!htbp]
\centering
\includegraphics[width=0.72\textwidth]{FB1_tracking.pdf}
\caption{Mode tracking across the thickness sweep, $L=\SI{200}{\micro\metre}$,
poling $y$, $48\times4$ mesh, $a_1\times a_2=3\times1$\,\si{\micro\metre},
$\mu=10^{-6}$\,\si{\coulomb\per\metre}. (a) Defect $1-\mathrm{MAC}$ between the
mode shape at each thickness and the shape at $h=\SI{10}{\micro\metre}$, on a
logarithmic scale, with the MAC itself printed above each bar; (b) frequencies
of the first two modes against thickness and their ratio, printed at each
point. The branch is separated from its neighbours over the whole range, which
makes the tracking unambiguous.}
\label{fig:modetrack}
\end{figure}

\paragraph{What the piezoelectric terms do on the bending device.} The bending
configuration is poled along $y$ with electrodes on the long faces, and the
crystal carries $e_{31}$ and $e_{33}$ as well, so the two effects are
separated inside one run. Table~\ref{tab:config} does it: three materials, one geometry, the same bending mode. With $\mu=0$ the
piezoelectric coupling of that mode is $1.0\times10^{-10}$ at every thickness,
because the piezoelectric charge $e_{31}\varepsilon_{xx}+e_{33}\varepsilon_{yy}$
is odd in $y-h/2$ and its integral over a face electrode vanishes; the device
is therefore not a piezoelectric device by design. With $e=0$ the flexoelectric
coupling runs from \num{1.76e-4} to \num{5.5e-3}, a factor $31$, while the full
crystal gives the $24$ reported above: the piezoelectric terms shift the
flexoelectric coupling by up to $24\,\%$ at the thinnest strip, through the
elastic stiffness of the same mode rather than through a charge they
contribute. The number quoted in the abstract and in Fig.~\ref{fig:thickness}
is the full-crystal one.

\begin{table}[!htbp]
\centering
\small
\caption{Piezoelectric and flexoelectric coupling on the \emph{same} bending
mode, one at a time and together: poling $y$, electrodes on the long faces,
$48\times4$ mesh, $L=\SI{200}{\micro\metre}$, window half-widths
$3\times\SI{1}{\micro\metre}$, $\mu=10^{-6}$\,\si{\coulomb\per\metre}. The
second column switches the piezoelectric constants off ($e=0$), the third
switches $\mu$ off, and the last pair reports the material of the paper. The
values of the fourth column are those plotted in
Fig.~\ref{fig:thickness}(a).}
\label{tab:config}
% table generated from the archived results
\begin{tabular}{@{}lcccc@{}}
\toprule
$h$ [\si{\micro\metre}] & $k^2$ piezo only & $k^2$ flexo only, $e=0$ & $k^2$ both, as reported & both/flexo $-1$\\
\midrule
10 & 1.0e-10 & 1.7569e-04 & 1.7437e-04 & -0.8\%\\
5 & 1.0e-10 & 6.5515e-04 & 6.5873e-04 & +0.5\%\\
2 & 1.0e-10 & 2.8454e-03 & 2.6416e-03 & -7.2\%\\
1 & 1.0e-10 & 5.4663e-03 & 4.1736e-03 & -23.6\%\\
\bottomrule
\end{tabular}

\end{table}

\paragraph{Numerical reproducibility of the thinnest point.} The
$h=\SI{1}{\micro\metre}$ point is the least reproducible number in the paper.
Its $k^2$ is a difference of two nearly equal eigenvalues of a pencil whose
condition number is $\sim10^{12}$ at that thickness, so it inherits the accuracy
of the eigensolver rather than of the model: four mathematically equivalent
eigensolver paths and four numpy/scipy stacks agree on $4.2\times10^{-3}$ and
on the factor $24$ to within $0.5\,\%$, and not more closely. The spread is the
range over its midpoint and is therefore a property of the model rather than of
the machine that computed it; at $h=\SI{10}{\micro\metre}$ the same spread is at
most $10^{-4}\,\%$. A better solver does not recover further digits, which is
why the two-significant-digit form is used here and in the abstract. The
platform-sensitivity table behind this statement is in
\ref{app:verification}.

\paragraph{The piezoelectric coupling is thickness independent.} Panel~(b) of
Fig.~\ref{fig:thickness} is the counterpart of panel~(a). The quasi-static
extensional coupling takes the values $0.55814$, $0.55843$, $0.55858$ and
$0.55861$ at $10$, $5$, $2$ and \SI{1}{\micro\metre}, that is it changes by
$0.09\,\%$ between \SI{10}{\micro\metre} and \SI{1}{\micro\metre}, against the
constitutive closed form $0.5608$ for the same material and the same
plane-stress kinematics. The piezoelectric effect is a bulk effect, so a thinner
strip is electromechanically neither better nor worse, whereas the flexoelectric
effect is a gradient effect and improves by a factor $24$ over the same range. A
device can therefore be thinned for the flexoelectric response without a penalty
in the piezoelectric one, and the crossover thickness is a property of the
window and of $\mu$, not of the piezoelectric constants.

\FloatBarrier
\subsection{Transient response and the memory model}
\label{sec:transient}

Figure~\ref{fig:ringdown} shows the ring-down of the extensional mode after an
initial tip displacement of \SI{10}{\nano\metre}, for eight Caputo orders
between $\alpha=1$ and $0.4$ at fixed $\tau=\SI{3}{\nano\second}$
($\omega\tau=0.13$): the tip displacement in panel~(a) and the open-circuit
voltage in panel~(b). The initial voltage, \SI{139}{\volt}, is set by the mode
shape and the electrode configuration and is independent of $\alpha$, as it
must be, since the memory term vanishes at $t=0^{+}$; the decay rate, by
contrast, depends on $\alpha$. The $\alpha=1$ curve is the exponential decay of
the classical Kelvin--Voigt element, while the fractional curves are not
exponential, and it is that difference which the damping map of
Fig.~\ref{fig:damping}(b) quantifies. The voltage is a linear functional of the
same modal coordinate and inherits the ring-down shape. Over the eight orders
plotted, the decrement of the tip displacement over the first two peaks,
converted with the exact relation $\zeta=\delta/\sqrt{4\pi^{2}+\delta^{2}}$ of
\eqref{eq:zeta}, runs from $0.063$ at $\alpha=1$ to a maximum of $0.104$ at
$\alpha=0.6$ and back to $0.085$ at $\alpha=0.4$
(Table~\ref{tab:ringdown}). For $\alpha<1$ it is a first-period decrement and
not a damping ratio, because those records are not exponential; a design that
takes the integer-order early decay as representative is in error by
$64\,\%$ at $\alpha=0.6$, and the sign of the error cannot be guessed from
$\alpha$ alone.

\begin{figure}[!htbp]
\centering
\includegraphics[width=\textwidth]{F9_ringdown.pdf}
\caption{Ring-down of the extensional mode of the
\SI{200}{\micro\metre}$\times$\SI{10}{\micro\metre} strip
($\SI{6.727}{MHz}$; $k^2=0.4583$ is the mode-based value
$(\omega_{oc}^{2}-\omega_{sc}^{2})/\omega_{oc}^{2}$ of that mode on the
$24\times2$ mesh of this figure, Table~\ref{tab:coupling}) for eight Caputo orders $\alpha=1$ to $0.4$:
(a) tip displacement, (b) open-circuit voltage. Modal memory applied to
$K_{\rm eff}$ with eight modes retained, $\tau=\SI{3}{\nano\second}$
($\omega\tau=0.13$, against $\tau=\SI{7.0974}{\nano\second}$ and
$\omega\tau=0.3$ in Table~\ref{tab:memory}), initial tip displacement
\SI{10}{\nano\metre}. Each curve is an independent L1 time march of $4000$
steps; all eight share the initial condition and differ only in $\alpha$. The
numbers quoted from these curves are the first-period decrements of
Table~\ref{tab:ringdown}, converted with \eqref{eq:zeta}; they are damping
ratios at $\alpha=1$ only.}
\label{fig:ringdown}
\end{figure}

\FloatBarrier
\subsection{Application: a \SI{200}{\micro\metre} by \SI{10}{\micro\metre}
piezoelectric--flexoelectric strip in open-circuit sensing}
\label{sec:application}

The worked case throughout is a \SI{200}{\micro\metre} long,
\SI{10}{\micro\metre} thick $\mathrm{BaTiO_3}$ strip clamped on one short edge,
poled along its length and electroded on the two ends. In this configuration
the strain in the active region is extensional rather than bending, so the
piezoelectric coupling is the larger of the two and the flexoelectric term is a
size-dependent correction rather than the main effect.
Table~\ref{tab:application} collects the numbers of the four
polarization/kinematics combinations of the same strip, each with the frequency
of the mode that carries its coupling. The fundamental of the extensional
configuration is at \SI{147.73}{kHz}, above the clamped-free value of a beam of
the same length because the gradient terms stiffen the member, while the
strongly coupled extensional mode sits at \SI{6725.8}{kHz} in plane stress.

\begin{table}[!htbp]
\centering
\small
\caption{The same \SI{200}{\micro\metre}$\times$\SI{10}{\micro\metre}
$\mathrm{BaTiO_3}$ strip in four polarization/kinematics combinations: the
mode-free coupling factor, the coupling of the strongest mode, the open-circuit
voltage at \SI{10}{\nano\metre} of tip motion, and the relative change in
coupling over $10\to1$\,\si{\micro\metre} of thickness. The last two rows are
the two electrical configurations of Fig.~\ref{fig:geometry}; they are not
additive.}
\label{tab:application}
% table generated from the archived results
\begin{tabular}{@{}llcccc@{}}
\toprule
poling & kinematics & mode & $f$ (kHz) & $k^2$ & $V$ at \SI{10}{\nano\metre} (V) \\
\midrule
y & plane strain & bending (4) & 6721.8 & 0.1749 & 3.48 \\
y & plane stress & bending (4) & 6361.0 & 0.1141 & 2.41 \\
x & plane strain & extensional (5) & 7161.0 & 0.4024 & 141.45 \\
x & plane stress & extensional (5) & 6725.8 & 0.4580 & 138.99 \\
\bottomrule
\end{tabular}

\end{table}

Three consequences follow from the runs and are worth stating explicitly.

\begin{enumerate}
\item \textbf{Polarization direction dominates everything else.} The
  length-extensional configuration with poling along the length gives
  $k^2=0.4580$ and \SI{138.99}{\volt} at \SI{10}{\nano\metre} of motion, while
  the bending configuration gives $0.1141$ and \SI{2.41}{\volt}: a factor $58$
  in voltage at equal amplitude and a factor four in $k^2$. A device that
  excites bending and senses through $e_{31}$ works against that handicap at
  this geometry, and the handicap is a property of the electrode
  arrangement rather than of $e_{31}$, since a symmetric strip bent between
  full through-thickness electrodes produces almost no net charge: the
  fundamental bending mode has $k^2=2.3\times10^{-10}$, the mode that couples
  through the thickness being the fourth.
\item \textbf{The quasi-static coupling of the same strip is $0.5581$} and is
  thickness independent (Sec.~\ref{sec:thickness}). This is the number to quote
  for a device operating far below its first resonance, and the number that
  remains meaningful when the strip is thinned for the flexoelectric
  contribution.
\item \textbf{The damping is set by $(\alpha,\tau)$ and by the memory matrix},
  not by a material constant. For the extensional mode of this strip at
  $\omega\tau=0.3$ and $\alpha=1$, $\zeta_{\rm eff}=0.15003$ with the memory
  on $K_{\rm eff}$ and $0.0675$ with it on $K_{uu}$, a ratio of $0.450$. At
  $\alpha=0.7$ the decay is not exponential and the quantity that survives is
  the first-period decrement: $0.194$ on $K_{\rm eff}$ against $0.083$ on
  $K_{uu}$, both first-period decrements and not damping ratios. A design that
  reports a single damping number without naming $\alpha$, $\tau$ and the
  memory matrix is under-specified.
\end{enumerate}

\FloatBarrier
